\documentclass[12pt]{article}
\usepackage{graphicx}
%\usepackage[a4paper,margin=1in]{geometry}
\usepackage{hyperref}
\usepackage{amsmath}
\usepackage{amssymb}
\usepackage{natbib} 
\bibliographystyle{plain}
\usepackage{siunitx}
\usepackage{bm}

\title{Conformations Search methods for molecular docking}
\author{Elvis Martis}
\begin{document}
\maketitle


%\chapter*{Conformational Search Methods}

\section{Basic concepts}
Understanding molecular conformation the three-dimensional arrangement of atoms in a molecule is fundamental to computational chemistry. A molecule's conformation determines its chemical and biological properties, including reactivity, solubility, and binding affinity to biological targets. In the context of this book, this conformational search is performed withing the binding site of receptor, which aims to locate the bioactive conformation. Unlike global or local minimum conformation in solution, bioactive conformtion is one of the low energy conformations that shows most appropriate shape complementary with the binding site. This maximises favourable interactions between the ligand and receptor atoms. The success of structure-based drug design methods depends on identifying this bioactive conformation. 

%verify the statements below
A typical small drug-like molecule may contain multiple rotatable bonds, each capable of adopting several energetically distinct rotameric states. With $n$ rotatable bonds and approximately 3 accessible rotameric states per bond, the conformational space contains roughly $3^n$ distinct conformations. For a molecule with 12 rotatable bonds typical for many pharmaceutically relevant compounds, this yields nearly 530,000 possible conformations. Manual identification becomes impractical, necessitating automated conformational search methods.


\subsection{Potential Energy Surface}

The potential energy of a molecule is expressed as a function of the nuclear/atomic coordinates\cite{CH2_Jensen2017,CH2_Cramer2013} given by:

\begin{equation}
  V(\mathbf{R}) = V(R_1, R_2, \dots, R_{3N}),
\end{equation}

where N is the number of atoms and $\mathbf{R}$ is the 3N-dimensional vector of Cartesian nuclear positions. For a non-linear molecule, the number of internal degrees of freedom is 3N - 6, after removal of overall translation and rotation. Formally, a potential energy surface (PES) can be defined as the hypersurface obtained by plotting $V(\mathbf{R})$ against a chosen set of internal coordinates (bond lengths, angles and torsions). Such high-dimensional PES cannot be visualised for human interpretation, therefore, we must restrict plotting only a few important internal coordinates. These important internal coordinates are also known as collective variables that might be enough explain the behaviour of the entire chemical or biochemical system.

\subsubsection{Stationary points and the Hessian matrix}
 While exploring the PES, there key parts of the PES that we are most interested in--locating the energy minimum. Mathematically, minimum or maximum points on the PES are located by differentiating the force field function describing the chemical system. In an ideal system, the minimum or maximum is defined as stationary points on the PES such that the first-order derivative of the function is zero and are classified by the eigenvalues of the Hessian matrix. The following equations illustrate this fact:
\begin{equation}
  \frac{\partial V}{\partial q_i} = 0
  \quad\text{for all internal coordinates } q_i,
\end{equation}

\begin{equation}
  H_{ij} = \frac{\partial^2 V}{\partial q_i\,\partial q_j}.
\end{equation}

A local minimum has all positive eigenvalues (no imaginary frequencies), whereas a first-order saddle point (transition state) has exactly one negative eigenvalue (one imaginary frequency). The topology of the PES around these points controls conformational equilibria and reaction rates\cite{CH2_Jensen2017}.

\subsubsection{Energy profiles and reaction coordinates}

In one dimension, for example rotation of a torsion angle, the PES reduces to an energy profile V(s) along a reaction coordinate s. Energy minima correspond to stable conformers, while barriers may correspond to transition structures. Such barriers also suggest 

\subsection{Internal coordinates}

For molecular mechanics representations, it is convenient to use internal coordinates such as:
\begin{itemize}
  \item Bond stretches r,
  \item Angle bends $\theta$,
  \item Improper torsions $\chi$ (for planarity/chirality),
  \item Proper torsions (dihedrals) $\phi$,
  \item Out-of-plane or improper torsion angles wherever necessary.
\end{itemize}

A generic functional form for the molecular mechanics potential energy is given by:

\begin{equation}
  V(\mathbf{q}) =
  \sum_\text{bonds} k_r (r - r_0)^2 +
  \sum_\text{angles} k_\theta (\theta - \theta_0)^2 +
  \sum_\text{torsions} V_\text{tor}(\phi) +
  V_\text{nonbonded},
\end{equation}

where $V_\text{nonbonded}$ computes van der Waals and electrostatic interactions, and $V_\text{tor}(\phi)$ is the torsional potential.

\subsection{Torsional potentials}

Torsional terms are usually represented as a Fourier series,

\begin{equation}
  V_\text{tor}(\phi) =
  \sum_{n} \frac{V_n}{2}\bigl[1 + \cos\bigl(n\phi - \gamma_n\bigr)\bigr],
\end{equation}

where $V_\text{n}$ are barrier heights, n the periodicity, and $\gamma_\text{n}$ phase offsets. The barrier height controls the rate of interconversion between conformers. Very low or negligible barrier makes the conformers indistiguishable, hence, considered equivalent for all practical purposes. 

If we consider simple organic molecules:
\begin{itemize}
  \item Ethane: $n = 3$, $\gamma = 0$, giving three equivalent staggered minima and three eclipsed maxima.
  \item Butane: interplay between $n=3$ and $n=1$ terms to distinguish anti- and gauche-staggered minima.
\end{itemize}


\subsection{Rotatable bonds}
Chemically, a rotatable bond can be defined as follows: 

\begin{itemize}
  \item A single (non-aromatic) bond ($\sigma$ bond) between heavy atoms
  \item That is not in a ring
  \item And is not terminal (i.e. not $X--CH_3$ or similar simple substituents)
  \item Excluding amide C-N bonds, and other formally conjugated single bonds with high torsional barriers.
\end{itemize}


\subsection{Conformational ensembles and thermodynamics}

Given a set of conformers i with relative energies $E_i$, the Boltzmann-weighted population at temperature T is

\begin{equation}
p_i = \frac{\exp(-E_i/k_\mathrm{B}T)}{\displaystyle\sum_j \exp(-E_j/k_\mathrm{B}T)},
\end{equation}

so that low-energy conformers dominate the bulk properties. Computational workflows therefore:

\begin{itemize}
  \item Sample a sufficiently diverse ensemble on the PES (system- and method-dependent).
  \item Reweight or rescore conformers with higher-level methods when needed.
  \item Consider entropic contributions associated with loss of conformational freedom upon binding.
\end{itemize}

\section{Metrics to ensure conformational sampling}

\subsection{Energy-Based Metrics}

Key metrics characterize conformational search results:

\begin{equation}
\Delta E = E_{\text{conformer}} - E_{\text{global minimum}}
\label{eq:relative_energy}
\end{equation}

The relative energy $\Delta E$ expresses each conformation's energy relative to the lowest-energy structure. The Boltzmann population at temperature $T$ predicts the probability of observing a particular conformation:

\begin{equation}
P_i = \frac{e^{-\Delta E_i / k_B T}}{\sum_j e^{-\Delta E_j / k_B T}}
\label{eq:boltzmann}
\end{equation}

where $k_B$ is Boltzmann's constant (1.987 cal·mol$^{-1}$·K$^{-1}$). At 298 K, conformations within 2 kcal/mol of the global minimum constitute $>10\%$ of the population and should be included in conformational ensembles.

The RMSD quantifies geometric similarity between two conformations:

\begin{equation}
\text{RMSD} = \sqrt{\frac{1}{N} \sum_{i=1}^{N} d_i^2}
\label{eq:rmsd}
\end{equation}

where $d_i$ represents the distance between atoms $i$ after optimal superposition. RMSD thresholds (typically 1.0--2.0 \r{A}) filter redundant conformations in search output.


\section{Conformational search methods}
One of the most challenging aspect of pose prediction is the conformational search algorithm. In general, there are two main classes of conformational search methods, systematic and stochastic.  The users do not always have a choice to select the search algorithm because the developers usually select based on their expertise and code development. Nevertheless, there are programs, such as AutoDock, that allow users to select the search algorithm. 

\subsection{Systematic Methods}
The systematic search algorithm rotates all possible rotatable bonds (single bonds in the ring structures cannot be rotated as freely as open-chain structures) systematically by a fixed interval to explore all possible conformations. Such a method can explore all possible conformations, however, the complexity increases exponentially as the number of rotatable bonds increases. There are pruning algorithms that act as a “bump check” that rule out torsion angles that lead to significant atomic overlaps. The deterministic nature of systematic sampling guarantees that the same input will yield the same set of conformers, which is advantageous for benchmarking and reproducibility. However, the combinatorial explosion of grid points with increasing rotatable bonds quickly becomes prohibitive, especially for highly flexible ligands, motivating the use of hybrid schemes or coarse‑grained grids in practice\cite{CH2_Beusen1996,CH2_Leach2001}.

\subsubsection{Grid-based enumeration} 
 
    The simplest systematic approach divides torsion angle space into discrete intervals (e.g., 10° or 30°) and evaluates all possible combinations. For a molecule with $n$ rotatable bonds and 36 intervals per 360° rotation:

\begin{equation}
N_{\text{conformations}} = 36^n
\label{eq:grid_enumeration}
\end{equation}

    This method guarantees finding the global minimum but becomes computationally prohibitive for large $n$ (combinatorial explosion). Grid enumeration works well for rigid molecules with 2--4 rotatable bonds but is impractical for flexible compounds.
    
    In the context of molecular docking, the search space is limited to the binding pocket, often provided as a grid. Therefore, the systematic conformational search engine operates by discretizing the relevant degrees of freedom typically torsional angles on a this grid.  It then generated and evaluates all combinations of torsion angles within that grid. For a molecule with n rotatable bonds, each torsion  is sampled at a fixed interval, producing a Cartesian product of values  $\phi_1 \in {\phi_1^{(1)}, \dots, \phi_1^{(k)}}, \phi_2 \in {\phi_2^{(1)}, \dots, \phi_2^{(k)}}$ and so on. At each grid point, the geometry is energy‑minimized, and the resulting conformers are clustered and ranked by energy or empirical scoring functions. Because the search is exhaustive within the chosen resolution, systematic methods can in principle locate all sterically allowed conformations of a given molecule, provided the grid is sufficiently fine and the force field adequately describes the intramolecular energetics\cite{CH2_Ferreira2015}.
  

\subsubsection{The Z method}
    The Z method\cite{CH2_Petrella2011} is a variant of grid-based conformational search protocols that employs statistical information on the rotamers. It partitions the conformational space into subspaces and then applies repeated, uniform changes to a common starting structure, generating large ensembles of trial conformations in parallel. This approach allows rigid systematic searches in high dimensional spaces by combining grid‑based sampling with statistical information from rotamer libraries. The Z Method can also be used hierarchically: low resolution searches identify approximate regions of interest, which are then refined with finer grids or energy minimization protocols. The key advantage is the ability to couple systematic search with probabilistic rotamer preferences, improving the sampling efficiency while preserving the exhaustiveness of the underlying grid. However, to the best of our knowlegde, this method has only been applied protein stucture refinement. We did not find any literature references indicating its use as conformational search tool for small molecules or in molecular docking. 

\subsection{Fragmentation method}

In this method, the molecule is broken down into individual fragments. The fragments are then docked in the most suitable sub-pocket in the binding site. The entire molecule is then sequentially built by adding the linkers. This follows a systematic search for the linker conformation that optimally holds the fragments in their respective sub-pockets. The fragmentation method selects the rigid molecular components, for example ring structures, as fragments, and the flexible components as linkers. In doing so, this method reduces the computational complexity compared to the systematic search by focusing on small components, more particularly the linkers. FlexX and DOCK are examples of the docking algorithms that employ this method as its conformational search algorithm.


\subsection{Stochastic Methods}
Stochastic techniques utilize random sampling and probabilistic methods to explore the conformational space. They work by randomly making changes, preferably, to the rotatable bonds and evaluates the energy following the change. Metropolis Monte Carlo, and Genetic algorithm are popular stochastic search methods that are used as conformational search tools in docking programs. 

\subsubsection{Monte Carlo} 

It works on the principle of random sampling for a predetermined number of iterations to locate optimum conformation\cite{CH2_Metropolis1953}. The search algorithm starts by making a random changes in any of randomly selected rotatable bond. Following which, it evaluates the energy and compares it with the starting conformation. If the energy of the new state is lower than the previous state, then the conformation is accepted. If the energy of the new state is higher than the previous state, then it assigns a random number between 0 and 1 and evaluates its probability against the Boltzmann distribution; if the random number assigned is greater than the probability, then the new state is accepted otherwise rejected. In case of the rejection, the algorithm restarts from the previous conformations. Glide docking program is also known to employ Monte Carlo simulation in its docking algorithm to improve its pose prediction accuracy.

In Metropolis Monte Carlo\cite{CH2_Metropolis1953}, a trial move from configuration \(\mathbf{x}\) to \(\mathbf{x}'\) with energy change \(\Delta E = E(\mathbf{x}') - E(\mathbf{x})\) is accepted with probability
\begin{equation}
    P_{\mathrm{acc}}(\Delta E;T) =
    \begin{cases}
        1, & \Delta E \le 0, \\
        \exp\!\left(-\dfrac{\Delta E}{k_{\mathrm{B}} T}\right), & \Delta E > 0.
    \end{cases}
\end{equation}
At high T, uphill moves are readily accepted, enabling broad exploration of conformational space; as T is lowered, the algorithm increasingly favors low-energy states, mimicking physical annealing \cite{CH2_Kirkpatrick1983}.

\subsubsection{Genetic Algorithm (GA)} 

Genetic algorithms are population-based metaheuristics inspired by biological evolution.\cite{CH2_Judson1993,CH2_Goodman1998} A GA iteratively evolves a population of conformers under the action of selection, crossover, and mutation, guided by a fitness function (typically related to energy or docking scores).

A basic GA cycle for conformational search is:

\begin{enumerate}
  \item Initialization: Generate an initial population of conformers (e.g., random torsion angles rotation).
  \item Evaluation: Compute a fitness value for each conformer using an energy function
  \item Selection: Select parent conformers with probability biased toward higher fitness (e.g., tournament selection, roulette wheel).
  \item Variation: Apply crossover and mutation operators in torsion space to generate offspring conformers.
  \item Local optimization: perform local geometry minimization to relax each offspring to the nearest local minimum.
  \item Replacement: Form the new population from parents and offspring (e.g., generational or steady-state replacement, elitism).
  \item Termination: Stop when convergence criteria are met (e.g., fixed number of generations, no improvement, or conformational coverage saturation).
\end{enumerate}


A critical design choice is how to encode conformations as chromosomes. For small molecules, a common representation is a vector of torsion angles: \cite{CH2_Judson1993,CH2_Goodman1998,CH2_Mekenyan1999}

\begin{itemize}
  \item Each gene corresponds to a rotatable single bond, double bond stereochemistry, or ring-pucker degree of freedom.
  \item The allele is a torsion angle (continuous, typically in $[-\pi, \pi)$), or a discrete bin for that torsion.
  \item Chirality and double-bond E/Z configuration can be handled as additional discrete variables if configuration is allowed to vary.
\end{itemize}
 

\subsection{Molecular Dynamics Simulations}

 Molecular dynamics (MD) simulates the atomic motions as a function of time by solving the Newton's equation of motion.  MD simulations and molecular docking are complementary\cite{CH2_guterres2020improving} that are commonly employed in structure-based drug discovery to understand the protein-ligand interactions and their dynamics aspects. For computational efficiency, molecular docking algorithms treat receptors atoms rigid and the ligand as flexible. This leads to incorrect pose prediction when induced fit binding is observed\cite{CH2_salmaso2018bridging}. Therefore, MD simulation is an elegant method to incorporate such induced fit effects. This can be achieved in two ways, as shown below:

\begin{itemize}
\item Pre-docking step to sample various conformations of the receptor without the influence of the ligand.  This can serve as multiple receptor conformations for docking. 
\item Post-docking step to refine the docked receptor-ligand conformation by allowing the complex to evolve to their realistic conformation, which would be more probable physiologically.
\end{itemize}

A full chapter is dedicated to MD and methods based on it for post-processing and enhancing the sampling of receptor-ligand dynamics (\textbf{Link MD chapter}). 

\subsection{Simulated Annealing}

Simulated annealing (SA) is a concept more prevalent in engineering science, where malleability in metals/alloys is employed to manufacture cables for a vast variety of applications. For example, a long steel cable can be manufactured from a block of steel by heating then beating and then cooling. This cycle is repeated until the desired shape, size and length is achieved, which dictates the tensile strength of the steel cable. This heating-beating-cooling-heating cycle is simulated annealing. The heating increases the softness of the steel by making the atoms move apart by gaining kinetic energy at high temperatures. Thus, making the beating process easier and preventing any fracture or permanent deformation. The cooling stage traps the atoms in the beaten state. At the end of each cycle, the shape of the steel block progressively approaches the desired shape and size without losing the inherent properties of steel.

Previously, discussed idea was implemented as a powerfull stochastic optimization technique\cite{CH2_Kirkpatrick1983} widely used in conformational search for molecular docking, where the goal is to locate low‑energy ligand and protein conformations compatible with experimental binding modes. The heating in this case, simply meant adding more velocity to the atoms so that the can surmount the nearest energy barrier. Similarly, the cooling step meant reducing the velocity such the system now falls in the nearest local minimum. The heating and cooling is carried out iterative to ensure maximum smapling of the PES. This mapping of thermodynamic annealing onto conformational space makes SA particularly suited to global‑optimization problems such as flexible‑ligand docking, where the search space is rugged and high‑dimensional\cite{CH2_Ghose1993,CH2_Naidoo1997,CH2_Corcho2000}.

The conceptual basis of SA is the Boltzmann distribution for a system with energy E at temperature \(T\). 
The equilibrium probability of a configuration with energy \(E\) is
\begin{equation}
    p(E;T) \propto \exp\!\left(-\frac{E}{k_{\mathrm{B}} T}\right),
\end{equation}
where \(k_{\mathrm{B}}\) is the Boltzmann constant. 



\begin{enumerate}
\item \textit{Heating phase}: Starting from an initial conformation, the system is rapidly heated to high temperature (typically 800--1000 K), providing energy to overcome conformational barriers.
\item \textit{Annealing phase}: Temperature is gradually reduced according to a cooling schedule, allowing the system to explore broader conformational space while progressively favoring lower-energy structures.
\item \textit{Minimization phase}: At low temperature (50--100 K), geometry optimization drives the system to a nearby local minimum.
\end{enumerate}

The SA protocol may be repeated several times from different initial conformations, generating diverse low energy structures. Effectiveness depends on cooling rate such that too rapid cooling traps the system in local minima, which too slow cooling demands excessive computation.

\subsection{Distance Geometry Methods}

Distance geometry (DG) provides an alternative formulation of conformational search in terms of interatomic distances rather than Cartesian coordinates.\cite{CH2_crippen_havel_book,CH2_havel_kuntz_crippen_1983,CH2_crippen_chem_dg_1991} The central idea is to work in distance space: define lower and upper bounds on selected interatomic distances, complete these to a full metric that is (approximately) realizable in 3D Euclidean space, then embed this metric back into Cartesian coordinates. This approach leads to efficient generation of diverse, approximately low-energy conformations using only geometric information and modest energetic refinement.\cite{CH2_spellmeyer_dg_conf_1997,CH2_riniker_landrum_2015}

\subsubsection{Molecular distance geometry}
The molecular distance geometry problem (MDGP) formalizes molecular conformation generation as follows:

Let $G = (V,E)$ be a graph whose vertices $V = \{1,\dots,n\}$ correspond to atoms. For each edge $(i,j) \in E$ we are given a distance interval

\begin{equation}
 \ell_{ij} \leq d_{ij} \leq u_{ij},
\end{equation}
where $d_{ij}$ is the (unknown) true distance between atoms $i$ and $j$, and $\ell_{ij}, u_{ij} \ge 0$ are lower and upper bounds derived from e.g.\ ideal bond lengths, angle geometry, van der Waals radii, NMR restraints, or crystallographic data.\cite{CH2_crippen_chem_dg_1991,CH2_havel_dg_review_1998}

Find coordinates $\bm{x}_1,\dots,\bm{x}_n \in \mathbb{R}^3$ such that
\begin{equation}
 \ell_{ij} \le \|\bm{x}_i - \bm{x}_j\| \le u_{ij}, \qquad \forall (i,j) \in E,
 \label{eq_mdgp}
\end{equation}
and specified chirality constraints for selected quadruples of atoms are satisfied.
For small molecules, E typically includes:
\begin{itemize}
  \item covalent bonds (tight bounds from ideal bond lengths),
  \item 1--3 and 1--4 distances derived from bond angles and torsion ranges,
  \item selected nonbonded distances (e.g. minimum separations to avoid clashes),
  \item optional experimental distances (e.g. NMR NOEs) when available.
\end{itemize}

A single molecule usually yields several solutions to \eqref{eq_mdgp}, corresponding to different conformers. Therefore, conformational search using distyance geomstry focuses on sampling the feasible region defined by the aforementioned constraints, typically followed by force field optimisation and clustering.\cite{CH2_spellmeyer_dg_conf_1997,CH2_riniker_landrum_2015}

\subsubsection{Classical Distance Geometry Algorithm for Conformer Generation}

The standard DG pipeline in chemical applications, sometimes referred to as the EMBED algorithm and its variants, consists of three major stages:\cite{CH2_havel_kuntz_crippen_1983,CH2_havel_dg_review_1998,CH2_spellmeyer_dg_conf_1997}
\begin{enumerate}
  \item \textit{Bound smoothing}: Given initial bounds on a sparse set of distances, bound smoothing propagates these constraints to non-neighboring atom pairs using metric inequalities. 
  The simplest are triangle inequalities: 
  \begin{equation}
  d_{ik} \le d_{ij} + d_{jk}, \qquad 
  d_{ik} \ge |d_{ij} - d_{jk}|
  \label{eq-Bound_smoothing}
  \end{equation}
  which lead to tightened upper and lower bounds $u_{ik}$ and $\ell_{ik}$ when $d_{ij}$, $d_{jk}$ are themselves bounded. Computationally, triangle-based smoothing method reduces to an all pairs shortest path problem on a suitably constructed weighted graph.This can be implemented in $\mathcal{O}(n^3)$ time or better for sparse molecular graphs.\cite{CH2_havel_kuntz_crippen_1983,CH2_havel_dg_review_1998}.
  

  \item \textit{Embedding}: After smoothing, one obtains approximate lower and upper bounds $\ell_{ij},u_{ij}$ for all $i,j$. A specific target distance matrix D is then sampled by choosing each $D_{ij}$ uniformly or normally within $[\ell_{ij},u_{ij}]$, subject to symmetry and zero diagonal. Embedding proceeds via the metric matrix, 
  \begin{equation}
  B = -\frac{1}{2} J D^{(2)} J
  \end{equation} 
  If B is exactly PSD with $\mathrm{rank}(B) \le 3$, then it corresponds to an exact Euclidean distance matrix (EDM) in $\mathbb{R}^3$ and the coordinates follow from the eigen-decomposition as above.
  
  \item \textit{Refinement}: Refinement improves the geometrical consistency and (optionally) the energy of the embedded structure. 
  A standard choice is to minimize a DG error function
\begin{equation}
 E(\{\bm{x}_i\}) =
 \sum_{(i,j)\in E} w_{ij}
 \left[
  f(d_{ij}(\{\bm{x}_i\}); \ell_{ij}, u_{ij})
 \right]^2,
\end{equation}
where $d_{ij}(\{\bm{x}_i\}) = \|\bm{x}_i - \bm{x}_j\|$ and $f$ penalizes deviations outside the allowed interval, e.g.
\begin{equation}
 f(d;\ell,u) =
 \begin{cases}
  d - u, & d > u,\\
  \ell - d, & d < \ell,\\
  0, & \text{otherwise}.
 \end{cases}
\end{equation}
Weights $w_{ij}$ can reflect confidence in the underlying constraints (for example, high for covalent geometry, lower for heuristic nonbonded bounds).\cite{CH2_crippen_chem_dg_1991,CH2_havel_dg_review_1998}

Optimization of E is typically performed using local methods such as conjugate gradients, quasi-Newton schemes, or simulated annealing. For small molecules, it is customary to combine DG error terms with a molecular mechanics force field:
\begin{equation}
 E_{\text{tot}} = E_{\text{DG}} + \lambda E_{\text{FF}},
\end{equation}
with $\lambda$ chosen to balance strict adherence to DG constraints against optimization of a realistic potential.\cite{CH2_spellmeyer_dg_conf_1997} This helps correct artifacts such as distorted aromatic rings or planar sp$^2$ centers that can arise in purely distance-based embeddings.\cite{CH2_riniker_landrum_2015}
\end{enumerate}


\subsection{Hybrid methods}

Stochastic and systematic search methods have their own pros and cons, and differ in their utility depending on the system under study. In this section, we discuss hybrid methods that use stochastic and systematic methods in tandem or in a cyclical manner to improve conformational search with better commputational efficiency than using either method standalone. 


\section{Practical Considerations}

\subsection{Conformational flexibility in acyclic systems}

For simple alkanes, conformational preferences are well rationalised by
\emph{Newman projections}:
\begin{itemize}
  \item Staggered conformers minimise torsional strain and often benefit from
        hyperconjugative stabilisation.
  \item Eclipsed conformers are higher in energy due to torsional and sometimes
        steric/electrostatic repulsions.
\end{itemize}
In butane, the anti-staggered conformer is lowest in energy, with gauche
conformers slightly higher but still populated at room temperature; this basic
model generalises to many drug-like substituents.

Hyperconjugation and stereoelectronic effects (e.g.the gauche effect,
anomeric effect) often bias torsional preferences away from purely steric
expectations. These effects are naturally captured in quantum chemical PES
scans and, when parametrised appropriately, in modern force fields.

\subsection{Conformational flexibility in ring systems}

\subsubsection{Ring strain components}

Cycloalkanes illustrate several types of strain:
\begin{itemize}
  \item \textit{Angle strain}: deviation of bond angles from the ideal tetrahedral value (\(\approx\SI{109.5}{\degree}\)).
  \item \textit{Torsional strain}: eclipsing interactions between neighbouring bonds.
  \item \textit{Steric (non-bonded) strain}: 1,3-diaxial interactions, transannular repulsions in medium/large rings, etc.
\end{itemize}
Small rings (cyclopropane, cyclobutane) are strongly angle- and torsionally strained; cyclopentane relieves torsional strain by puckering; cyclohexane adopts near strain-free conformations.

\subsubsection{Cyclohexane: chair, boat and twist-boat}

In cyclohexane, the most important conformations are:
\begin{itemize}
  \item Two equivalent chair conformers, related by a ring flip.
  \item Higher-energy boat and twist-boat conformations.
\end{itemize}

In the ideal chair, all C--C bonds are staggered and the internal bond angles are close to tetrahedral. A ring flip exchanges axial and equatorial positions; for monosubstituted cyclohexanes, the equilibrium strongly favours the equatorial conformer for bulky substituents due to reduced 1,3-diaxial repulsions.

From a computational perspective, the PES in the relevant ring-puckering coordinates shows:
\begin{itemize}
  \item Two deep minima (chairs),
  \item A pathway via twist-boat and boat conformations connecting them,
  \item Barriers that set the timescale for ring inversions.
\end{itemize}

\subsubsection{Conformational analysis of substituted and heterocyclic rings}

Substitution and heteroatoms modify ring conformations and relative energies:

\begin{itemize}
  \item Electronegative or $\sigma$-accepting substituents can introduce stereoelectronic effects (e.g.\ anomeric effect in pyranoses).
  \item Heteroatoms change bond lengths/angles and steric profiles, altering A-values (axial penalty) and flipping preferences.
  \item Fused ring systems transmit conformational bias conformational transmission, important in polycyclic scaffolds used in medicinal chemistry.
\end{itemize}
Accurate treatment of these systems often requires a combination of high-quality force fields and validation against quantum chemical PES scans or experiment.

\subsection{Accounting for solvation effects}

Bioactive conformations depend on the solvent environment. Implicit solvation models (Generalized Born, Poisson-Boltzmann) add solvation energy terms, affecting the conformational landscape relative to gas-phase calculations. 

\subsection{Clustering of Results}

Why we need clustering? and how does this ensure use diverse conformations for next step.

RMSD:

\[
\text{RMSD} = \sqrt{\frac{1}{N} \sum_i |\mathbf{r}_i - \mathbf{r}_i'|^2}
\]

Clusters identify recurring poses.

\subsection{Validation Metrics}

Conformational search quality is evaluated by:

\begin{enumerate}
\item \textit{Recovery of Crystallographic Conformations}: For validation datasets with known bioactive structures (e.g., protein-bound ligand conformations), the minimum RMSD from the ensemble to experimental structure indicates success. Recovery within 2 \r{A} RMSD is considered good.

\item \textit{Conformational Diversity}: The range of geometries in the final ensemble should comprehensively represent the conformational space. Metrics include distribution of dihedral angles and range of relative energies.

\item \textit{Computational Efficiency}: Wall-clock time and peak memory usage relative to molecular size determine practical viability. For drug discovery pipelines screening millions of compounds, efficiency is paramount.
\end{enumerate}

\section{Summary}
Conformational search is the computational engine that enables docking. Understanding both the physics of the scoring function and the algorithms used to explore conformational space is essential for reliable predictions. Hybrid strategies that combine stochastic global exploration with local refinement currently offer the best trade-off between efficiency and accuracy.



\bibliography{conform_search}

\end{document}